\documentclass[runningheads,a4paper]{llncs}

\usepackage{amsmath,amssymb,amsfonts}
\usepackage{bm}
\usepackage{graphicx}
\usepackage{booktabs}
\usepackage{multirow}
\usepackage{xcolor}
\usepackage{subcaption}
\usepackage{algorithm}
\usepackage{algpseudocode}
\usepackage{tikz}
\usetikzlibrary{arrows.meta, positioning, shapes.geometric, fit}
\usepackage{hyperref}
\usepackage{luatexja}
\usepackage[match]{luatexja-fontspec}
\setmainjfont{Hiragino Mincho ProN}
\setsansjfont{Hiragino Kaku Gothic ProN}

% llncs は theorem/lemma/definition/remark/example/problem を提供済み
% 日本語名称に上書き
\renewcommand{\definitionname}{定義}
\renewcommand{\lemmaname}{補題}
\renewcommand{\remarkname}{注意}
\renewcommand{\problemname}{問題}
\renewcommand{\examplename}{例}

\begin{document}

\title{$d$ 次元 Apriori-Window: 多次元トランザクションデータにおける\\
密集アイテムセット領域の厳密抽出}

\author{高安 奏多\inst{1} \and 高橋 泰平\inst{1} \and 栗原 聡\inst{2} \and 須賀 聖\inst{3}}
\authorrunning{高安 奏多 ら}
\titlerunning{3次元 Apriori-Window: 密集アイテムセット領域の厳密抽出}

\institute{
  慶應義塾大学 理工学研究科，日本\\
  \email{kanata1222@keio.jp}
  \and
  慶應義塾大学 理工学部，日本
  \and
  関西大学 ビジネスデータサイエンス学部，日本
}

\maketitle

\begin{abstract}
Apriori-Window は時刻付きトランザクションデータに対してスライディングウィンドウ密集区間を
直接かつ厳密に発見する手法であるが、1次元の時間軸のみを扱うため、
空間的に局在するパターンを識別できない問題がある。
本稿では Apriori-Window を
任意の $d$ 次元グリッド（時間 + $d{-}1$ 個の空間・カテゴリ次元）に拡張した
\textbf{$d$ 次元 Apriori-Window アルゴリズム}を提案する。
3つの高速化技術のうち $t$ 軸ストライドと Apriori 枝刈りは前身手法の直接拡張であり、
$d$ 次元接頭辞和（$2^d$ コーナー包除）は前身手法の二分探索（$O(\log n)$）を
$O(2^d)$ 定数時間クエリに改善する新技術である。これらを組み合わせ、
計算量をナイーブ実装の $O\!\left(M \cdot G \cdot \prod W_k\right)$ から
$O\!\left(M \cdot G + N \cdot \prod_{k\geq1} W_k\right)$ に削減する。

合成データ実験（$d{=}3$）では、
提案手法が植え付けパターン3件を Precision $= 1.0$・Recall $= 1.0$ で完全回収する一方、
1次元相当設定は93〜100件の偽陽性パターンを生成した（Precision $\approx 0.03$）。
次元数スケーラビリティ実験（$d{=}1$〜$5$）では実行時間の次元依存性を定量化した。
$N = 100{,}000$ トランザクション・$d{=}3$ で約 1.4 秒（Rust + Rayon 並列化）を達成した。
\end{abstract}

\keywords{頻出パターンマイニング, 多次元データ, スライディングウィンドウ,
密集領域検出, Apriori, $d$ 次元接頭辞和}

% ==============================================================
\section{はじめに}\label{sec:intro}
% ==============================================================

頻出パターンマイニング~\cite{agrawal1994apriori,han2000mining}は
トランザクションデータ中に繰り返し出現するアイテムセットを発見する基盤技術であり、
小売分析、医療データ解析、Web ログ分析など幅広い分野に応用されている。

Apriori-Window~\cite{aprioriwindow}は、
タイムスタンプ付きトランザクションデータベースに対して、
スライディングウィンドウのサポートが閾値 $\theta$ 以上となる
\emph{密集区間}とそこで頻出するパターンを効率的に発見するアルゴリズムである。
同手法は $t$ 軸ストライド最適化と Apriori 枝刈りを組み合わせることで、
候補区間の評価コストを大幅に削減し、
LPFIM や LPPM などの既存手法を性能で上回ることを示している。

しかし現実のトランザクションデータは時刻だけでなく\textbf{空間的位置情報}
（店舗座標、地域コード、GPS 位置等）を伴うことが多い。
例えば、小売チェーンでは各トランザクションが特定の店舗（座標）で発生し、
疫学データでは症例の発生場所が記録される。
このような\emph{時空間トランザクションデータ}において、
「沿岸部店舗のみで集中的に売れるアイテムセット」のような
\textbf{時空間局所的な密集パターン}を発見することは、
時間軸のみを扱う Apriori-Window では不可能である。
全店舗データを時間軸に集約してしまうと、
空間的に局在するパターンが背景ノイズに埋もれてしまうからである。

\subsection{技術的課題}

Apriori-Window を1次元（時間軸）から
任意の $d$ 次元グリッド（時間 + $d{-}1$ 個の空間・カテゴリ次元）に
拡張するには以下の課題がある。

\begin{enumerate}
\item \textbf{支持曲面の計算爆発}：
$d$ 次元グリッド上の全パターン・全グリッド点でサポートを
評価する単純実装は $O(M \cdot G \cdot \prod_{k=0}^{d-1} W_k)$ の計算量を要する。
$G = \prod N_k$ はグリッドサイズ、
$M$ はパターン数、$\prod W_k$ はウィンドウ体積である。

\item \textbf{次元増加による候補爆発}：
空間次元が加わるごとにグリッドサイズ $G$ が指数的に増大し、
各パターンについてすべての点を評価するコストが膨大になる。

\item \textbf{$t$ 軸ストライドの $d$ 次元化}：
Apriori-Window の1次元ストライドアルゴリズム~\cite{aprioriwindow}を
各空間ウィンドウ $(s_1, \ldots, s_{d-1})$ に独立に適用できるよう構造化する必要がある。

\item \textbf{密集領域定義の一般化}：
1次元の「密集区間」を $d$ 次元の「密集領域（$2d$-近傍連結成分）」に自然に拡張し、
Apriori 枝刈りを密集セル集合の交差として定式化する必要がある。
\end{enumerate}

\subsection{本研究の貢献}

本稿では Apriori-Window~\cite{aprioriwindow}を
任意の $d$ 次元グリッド（時間 + $d{-}1$ 個の空間・カテゴリ次元）に拡張した
\textbf{$d$ 次元 Apriori-Window アルゴリズム}を提案する（\ref{sec:method}節）。
貢献は以下の3点である。

\begin{enumerate}
\item \textbf{$d$ 次元接頭辞和による $O(2^d)$ ウィンドウクエリ}：
出現グリッド $I_P$ に対して $\prod (N_k{+}1)$ の接頭辞和を構築し、
$2^d$ コーナー包除原理により任意ウィンドウのサポートを定数時間で取得する。
これにより接頭辞和構築 $O(G)$ と合わせて、
ナイーブな $O(G \cdot \prod W_k)$ サポート計算を $O(G)$ に削減する。

\item \textbf{$t$ 軸ストライドの $d$ 次元拡張}（$k = 1$）：
Apriori-Window の $t$ 軸ストライドアルゴリズム~\cite{aprioriwindow}を
各空間ウィンドウ $(s_1, \ldots, s_{d-1})$ の空間イベントリストに独立に適用することで、
不要な時間スキャンを排除する（補題~\ref{lem:stride}で正確性を証明）。

\item \textbf{$d$ 次元 Apriori 枝刈り}（$k \geq 2$）：
Apriori-Window の候補区間交差枝刈り~\cite{aprioriwindow}を
$d$ 次元密集セル集合の交差として拡張する（補題~\ref{lem:apriori}）。
交差集合が空のパターンは $I_P$ 構築を省略でき、大幅な枝刈りが可能である。
\end{enumerate}

合成データ実験（\ref{sec:experiments}節）では、
$d{=}3$ 設定が植え付け3パターンを Precision $= 1.0$、
Recall $= 1.0$ で完全回収する一方、
1次元相当設定（空間ウィンドウ全体をカバー）では背景ノイズ由来の偽陽性パターンが
多数検出される（Precision $\approx 0.03$）ことを示し、
空間次元の明示的なモデル化の重要性を実証する。
また $N = 100{,}000$ トランザクション・$d{=}3$ で約 1.4 秒（Rust + Rayon 並列化）を達成し、
次元数スケーラビリティ（$d{=}2{\sim}5$）も定量化した。

% ==============================================================
\section{関連研究}\label{sec:related}
% ==============================================================

本研究は、時空間トランザクションデータから密集アイテムセット領域を発見する問題を扱う。
この問題は複数の研究領域にまたがるため、
各領域の代表的手法と本研究の相違点を整理する。

\subsection{頻出パターンマイニングとAprioriアルゴリズム}

Agrawal ら~\cite{agrawal1994apriori}のAprioriアルゴリズムは、
アイテムセットの単調性（部分集合は常に親集合以上のサポートを持つ）を利用した
枝刈りにより、候補数を大幅に削減する。
Han ら~\cite{han2000mining}のFP-Growthは接頭辞木を用いてさらに高速化する。
これらは静的なトランザクションデータベースを対象とし、
時間的・空間的なパターン分布の変化は扱わない。

本研究は Apriori の枝刈り原理を $d$ 次元密集セル集合に拡張し、
候補セルの交差集合が空の場合に $I_P$ 構築を省略する（\S\ref{sec:alg_k2}）。

\subsection{Apriori-Window（直接の前身研究）}

Apriori-Window~\cite{aprioriwindow}は、
タイムスタンプ付きトランザクションデータに対して
スライディングウィンドウサポートが閾値を超える密集区間を
直接かつ厳密に発見するアルゴリズムである。
$t$ 軸ストライドにより非密集区間を効率的にスキップし、
Apriori 枝刈りにより $(k{-}1)$-部分集合に密集区間がないパターンを早期除去する。
比較実験では LPFIM~\cite{han2000mining}、LPPM、Recur を大幅に上回る性能を示した。

本研究はこの手法を1次元（時間軸のみ）から
任意の $d$ 次元（時間 + $d{-}1$ 個の空間・カテゴリ次元）に拡張する。
全空間次元のウィンドウ幅を1とすると本研究の問題設定・アルゴリズムは
Apriori-Window に帰着する（注意~\ref{rem:1d_reduction}、補題~\ref{lem:stride}）。

\subsection{時空間パターンマイニング}

軌跡パターンマイニング~\cite{giannotti2007trajectory}は
移動オブジェクトの軌跡から時空間パターンを発見する。
空間コロケーションマイニング~\cite{shekhar2001colocation}は
空間的に共起する特徴を発見する。
これらは空間的分布を扱うが、
スライディングウィンドウによるサポートの時空間集計や
密集領域の連結成分抽出は対象としない。

\subsection{密部分テンソルマイニング}

DenseAlert~\cite{shin2017densealert}やM-Zoom~\cite{shin2018mzoom}は
多次元テンソルにおける密ブロックを検出する。
これらはテンソル全体のランク分解に基づく異常検出が主目的であり、
アイテムセットパターンのスライディングウィンドウサポートや
$2d$-近傍連結成分としての密集領域抽出は扱わない。

\subsection{空間スキャン統計}

Kulldorff~\cite{kulldorff1997spatial}の空間スキャン統計は
地理的領域における疾患クラスターを検出する代表的手法である。
矩形窓を用いた時空間拡張~\cite{kulldorff2001prospective}も提案されているが、
アイテムセットパターンのサポートではなく点過程の密度を対象とし、
本研究とは問題設定が異なる。

\subsection{本研究の位置づけ}

表~\ref{tab:related}に各手法との比較を示す。
本研究は、Apriori-Window~\cite{aprioriwindow}の $t$ 軸ストライドと Apriori 枝刈りを
$d$ 次元接頭辞和・空間イベントリストと組み合わせることで、
多次元トランザクションデータから密集アイテムセット領域を効率的に発見する唯一のフレームワークである。

\begin{table}[t]
\centering
\caption{関連手法との比較。
$\checkmark$: 対応、$\times$: 非対応。}
\label{tab:related}
\small
\begin{tabular}{lccc}
\toprule
手法 & アイテムセット & 時空間密集領域 & スライディングウィンドウ \\
\midrule
Apriori~\cite{agrawal1994apriori}               & $\checkmark$ & $\times$     & $\times$     \\
LPFIM / LPPM / Recur                            & $\checkmark$ & $\times$     & $\checkmark$ \\
Apriori-Window~\cite{aprioriwindow}             & $\checkmark$ & $\times$ (1D) & $\checkmark$ \\
空間スキャン統計~\cite{kulldorff1997spatial}     & $\times$     & $\checkmark$ & $\times$     \\
DenseAlert~\cite{shin2017densealert}             & $\times$     & $\checkmark$ & $\times$     \\
コロケーション~\cite{shekhar2001colocation}      & $\checkmark$ & $\checkmark$ & $\times$     \\
\midrule
\textbf{本研究}                                  & $\checkmark$ & $\checkmark$ ($d$次元) & $\checkmark$ \\
\bottomrule
\end{tabular}
\end{table}

% ==============================================================
\section{問題定義}\label{sec:problem}
% ==============================================================

本節では、まず Apriori-Window~\cite{aprioriwindow}が扱う1次元設定を
整理し（\ref{sec:1d_review}節）、
続いてそれを $d$ 次元に拡張した本研究の問題定義を与える（\ref{sec:nd_extension}節）。

% ==============================================================
\subsection{Apriori-Window（1次元）の復習}\label{sec:1d_review}
% ==============================================================

$\mathcal{I} = \{i_1, \ldots, i_m\}$ をアイテム集合、
$\mathcal{D}_{\mathrm{1D}} = \{(t_n, d_n)\}_{n=1}^N$ を
タイムスタンプ付きトランザクション列とする。
$t_n \in \mathbb{Z}_{\geq 0}$ は整数タイムスタンプ（または離散化済みビン）、
$d_n \subseteq \mathcal{I}$ はアイテム集合である。

パターン $P \subseteq \mathcal{I}$ の\emph{出現集合}を
$\mathrm{occ}(P) = \{t_n : P \subseteq d_n\}$ と定義する。

\begin{definition}[区間サポート~\cite{aprioriwindow}]\label{def:sup1d}
ウィンドウ幅 $W_t$ のもとで、ウィンドウ左端 $l$ に対する
パターン $P$ の\emph{区間サポート}を
\[
  \mathrm{sup}(P,\, l,\, W_t)
  = \bigl|\{t \in \mathrm{occ}(P) : l \leq t < l + W_t\}\bigr|
\]
と定義する。
\end{definition}

\begin{definition}[密集窓・密集区間~\cite{aprioriwindow}]\label{def:dense1d}
サポート閾値 $\theta$ に対し、
$\mathrm{sup}(P, l, W_t) \geq \theta$ を満たすとき、
窓 $[l, l+W_t]$ を $P$ の\emph{密集窓}と呼ぶ。
密集窓の開始点 $l$ が連続する最長区間 $[s, e]$ を
$P$ の\emph{密集区間}と呼ぶ。
\end{definition}

Apriori-Window~\cite{aprioriwindow}は、
$t$ 軸ストライドによる不要スキャンの排除と
候補区間の交差に基づく Apriori 枝刈りにより、
全パターンの密集区間を効率的に列挙する。

% ==============================================================
\subsection{$d$ 次元への拡張}\label{sec:nd_extension}
% ==============================================================

本研究では、各トランザクションが時刻 $c_0$（時間軸）に加えて
$d-1$ 個の追加座標 $(c_1, \ldots, c_{d-1})$（空間・カテゴリなど）を持つ
\emph{$d$ 次元トランザクションデータ}を対象とする。

\paragraph{記法.}
$\mathcal{D} = \{(\bm{c}_n, d_n)\}_{n=1}^N$ を
\emph{$d$ 次元トランザクション列}とする。
ここで $\bm{c}_n = (c_n^{(0)}, c_n^{(1)}, \ldots, c_n^{(d-1)}) \in \mathbb{R}^d$ は
座標ベクトル（第0軸は時間）、
$d_n \subseteq \mathcal{I}$ はアイテム集合である。

\begin{remark}[データモデルの前提]\label{rem:one_per_cell}
理論的には、各トランザクションは $d$ 次元空間の1点に対応する。
すなわちビン化後の各格子点には高々1件のトランザクションが存在し、
出現グリッド $I_P$ は指示関数（$\{0,1\}$ 値）として機能する。
実応用では、例えば犯罪記録データ（ChicagoCrime）や小売 POS データにおいて、
同一時刻・同一場所で複数のイベントが発生することは稀である。
十分粗いビン幅 $\bm{\Delta}$ を選ぶことで、同一ビンへの複数イベントの重複を回避できる。
ビン内に複数のイベントが集約される場合は、
そのビンに含まれる全イベントのアイテム集合の和集合を
当該ビンの単一トランザクションとして扱えばよい。
\end{remark}

\begin{definition}[$d$ 次元ビンインデックス]\label{def:bin_nd}
ビン幅ベクトル $\bm{\Delta} = (\Delta_0, \ldots, \Delta_{d-1})$ のもとで、
トランザクション $(\bm{c}, d)$ のビンインデックスを
\[
  \bm{b}(\bm{c}) = \bigl(\lfloor c^{(k)} / \Delta_k \rfloor\bigr)_{k=0}^{d-1}
\]
と定義する。グリッドサイズを $G = \prod_{k=0}^{d-1} N_k$ とする
（$N_k$ は第 $k$ 次元のビン数）。
\end{definition}

\begin{definition}[$d$ 次元アイテム出現グリッド]\label{def:item_grid_nd}
パターン $P \subseteq \mathcal{I}$ の\emph{出現グリッド} $I_P$ を
\[
  I_P[\bm{b}]
  = \bigl|\{n : \bm{b}(\bm{c}_n) = \bm{b},\; P \subseteq d_n\}\bigr|
\]
と定義する。注意~\ref{rem:one_per_cell}の前提のもとでは
$I_P[\bm{b}] \in \{0, 1\}$（指示関数）であり、
ビン $\bm{b}$ にパターン $P$ を含むトランザクションが存在するか否かを表す。
\end{definition}

\begin{definition}[$d$ 次元スライディングウィンドウサポート]\label{def:support_nd}
ウィンドウビン幅ベクトル $\bm{W} = (W_0, W_1, \ldots, W_{d-1})$ のもとで、
グリッド点 $\bm{s} = (s_0, s_1, \ldots, s_{d-1})$ を開始点とする
ウィンドウ内のパターン $P$ の\emph{サポート}を
\begin{equation}\label{eq:support_nd}
  S_P(\bm{s})
  = \sum_{\bm{b} \in \prod_{k=0}^{d-1} [s_k,\, s_k + W_k)}
    I_P[\bm{b}]
\end{equation}
と定義する。全グリッド点 $\bm{s}$ にわたる $S_P(\bm{s})$ を
$P$ の\emph{サポートサーフェス}と呼ぶ。
\end{definition}

\begin{remark}[1次元設定への帰着]\label{rem:1d_reduction}
$d = 1$（時間軸のみ）のとき、$S_P(s_0) = \mathrm{sup}(P, s_0, W_0)$
（定義~\ref{def:sup1d}）に帰着する。
すなわち本研究の問題設定は Apriori-Window~\cite{aprioriwindow}の真の拡張である。
さらに $d = 3$（時間 $\times$ 空間 $x \times$ 空間 $y$）のとき、
本稿で示す3次元実験結果が得られる。
\end{remark}

\begin{definition}[密集セル]\label{def:dense_cell_nd}
サポート閾値 $\theta$ に対し、
グリッド点 $\bm{s}$ が
$S_P(\bm{s}) \geq \theta$ を満たすとき、
$\bm{s}$ を\emph{密集セル}と呼ぶ。
パターン $P$ の密集セル集合を
$\mathcal{L}_\theta(S_P)$ と表す。
\end{definition}

\begin{definition}[密集領域]\label{def:dense_region_nd}
密集セル集合 $\mathcal{L}_\theta(S_P)$ の
$2d$-近傍連結成分（各次元の正負方向 $\pm \bm{e}_k$ を隣接と定義）の
各成分を\emph{密集領域}と呼ぶ。
$d = 3$ のとき $2d = 6$ 近傍（辺共有隣接）となる。
\end{definition}

1次元の「密集区間」（定義~\ref{def:dense1d}）は、
$d = 1$ のときの密集領域に対応する。

\paragraph{Apriori 包含定理.}

\begin{lemma}[密集領域の単調性]\label{lem:apriori}
任意のパターン $P$ とその部分集合 $Q \subset P$ に対して、
$\mathcal{L}_\theta(S_P) \subseteq \mathcal{L}_\theta(S_Q)$ が成立する。
\end{lemma}

\begin{proof}
任意のビン $\bm{b}$ に対し、$P \subseteq d_n$ ならば $Q \subseteq d_n$ であるから
$I_P[\bm{b}] \leq I_Q[\bm{b}]$。
接頭辞和の単調性より $S_P(\bm{s}) \leq S_Q(\bm{s})$。
\end{proof}

補題~\ref{lem:apriori}は Apriori-Window の候補区間交差枝刈り~\cite{aprioriwindow}の
$d$ 次元版であり、$k \geq 2$ の枝刈りの根拠である（\ref{sec:alg_k2}節）。

\subsection{$d$ 次元密集領域検出問題}

\begin{problem}[$d$ 次元密集領域検出]\label{prob:dense_detection}
\textbf{入力}：$d$ 次元トランザクション列 $\mathcal{D}$、
ビン幅 $\bm{\Delta} \in \mathbb{R}^d$、
ウィンドウビン数 $\bm{W} \in \mathbb{Z}_{>0}^d$、
サポート閾値 $\theta$、最大パターン長 $L_{\max}$。

\textbf{出力}：
各パターン $P$ に対し、$P$ の密集領域の全集合
$\{\mathcal{R}_1, \mathcal{R}_2, \ldots\}$（各 $\mathcal{R}_j$ は密集セルの連結成分）。
\end{problem}

問題~\ref{prob:dense_detection}の計算量は
サポートサーフェス $S_P$ の計算コストに支配される。
単純実装の $O\!\left(M \cdot G \cdot \prod_{k=0}^{d-1} W_k\right)$ を
提案アルゴリズム（\ref{sec:method}節）では
$O\!\left(M \cdot G + N \cdot \prod_{k=1}^{d-1} W_k\right)$ に削減する
（第0軸を時間軸 $t$ とする）。

% ==============================================================
\section{提案手法: $d$ 次元 Apriori-Window アルゴリズム}
\label{sec:method}
% ==============================================================

問題~\ref{prob:dense_detection}を解くための
\textbf{$d$ 次元 Apriori-Window アルゴリズム}を提案する。
以下の3つの高速化技術を組み合わせ（Algorithm~\ref{alg:md_apriori}）、
計算量を
$O\!\left(M \cdot G + N \cdot \prod_{k=1}^{d-1} W_k\right)$ に抑える
（第0軸を時間軸 $t$ とする）。
各技術は Apriori-Window~\cite{aprioriwindow}の
対応する1次元技術の直接拡張である（表~\ref{tab:correspondence}）。

\begin{table}[t]
\centering
\caption{Apriori-Window~\cite{aprioriwindow}と本研究の技術対応表。}
\label{tab:correspondence}
\small
\begin{tabular}{ll}
\toprule
Apriori-Window（1次元）& 本研究（$d$ 次元）\\
\midrule
二分探索 $O(\log n)$ 区間クエリ & $d$ 次元接頭辞和 $O(2^d)$ ウィンドウクエリ（新規）\\
$t$ 軸ストライド（$k=1$） & 空間イベントリスト + $t$ 軸ストライド（$k=1$）\\
候補区間交差による枝刈り & 候補セル集合交差による枝刈り（$k\geq2$）\\
密集区間（1次元連結成分） & 密集領域（$2d$-近傍連結成分）\\
\bottomrule
\end{tabular}
\end{table}

% ==============================================================
\subsection{$d$ 次元接頭辞和による $O(2^d)$ ウィンドウクエリ}\label{sec:prefix_sum}
% ==============================================================

出現グリッド $I_P$（サイズ $\prod N_k$）に対し、
各次元を $+1$ パディングした $d$ 次元接頭辞和 $\mathrm{PS}_P$（サイズ $\prod (N_k+1)$）を構築する：
\begin{equation}\label{eq:prefix_sum}
  \mathrm{PS}_P[\bm{b}]
  = \sum_{\bm{b}' \prec \bm{b}} I_P[\bm{b}'],
\end{equation}
ここで $\bm{b}' \prec \bm{b}$ は成分ごとに $b'_k < b_k$（全次元）を意味する。
$\mathrm{PS}_P$ は各次元を順番に累積和をとることで $O(G)$ で構築できる。
構築後、任意のウィンドウサポートは $2^d$ コーナーの包除原理により
\begin{equation}\label{eq:query}
  S_P(\bm{s})
  = \sum_{\bm{c} \in \{0,1\}^d}
    (-1)^{d - |\bm{c}|_1}
    \mathrm{PS}_P[\bm{s} + \bm{c} \odot \bm{W}]
\end{equation}
で $O(2^d)$ に取得できる（$\bm{c} \odot \bm{W}$ は成分積）。
$k \geq 2$ ではこれを全グリッド点に適用し、
サポートサーフェス $S_P$ を $O(2^d \cdot G)$ で計算する。
実用的な $d$（$\leq 5$ 程度）では $2^d$ は定数とみなせる。

% ==============================================================
\subsection{$t$ 軸 1 次元ストライド最適化（$k = 1$）}\label{sec:t_stride}
% ==============================================================

$k = 1$（単体アイテム）の密集セル検出では、
Apriori-Window の $t$ 軸ストライド~\cite{aprioriwindow}を
各空間ウィンドウ $(s_1, \ldots, s_{d-1})$ に独立に適用する。

\paragraph{空間イベントリストの構築.}
トランザクション $(\bm{c}_n, d_n)$（ビン化後: $\bm{b}_n$）がアイテム $i$ を含む場合、
$i$ が関係する空間ウィンドウ開始点の第 $k$ 成分（$k \geq 1$）の範囲は
\[
  s_k \in [\max(0,\, b_n^{(k)} - W_k + 1),\; b_n^{(k)}]
\]
であるから、全空間次元の組み合わせ $(s_1, \ldots, s_{d-1})$ に
$b_n^{(0)}$（時間ビン）を追加する。
構築コストは $O\!\left(N_i \cdot \prod_{k=1}^{d-1} W_k\right)$。
$d = 1$（時間のみ）のとき単一リスト $L_i() = \mathrm{occ}(\{i\})$ に帰着する。

\paragraph{$t$ 軸ストライド.}
各 $(s_1, \ldots, s_{d-1})$ について、ソート済みリスト $L_i(s_1, \ldots, s_{d-1})$ に対し
ウィンドウ左端 $l$ を動かして密集窓の開始点を求める
（Algorithm~\ref{alg:t_stride}）。
本アルゴリズムは Apriori-Window の $t$ 軸ストライドと同一の
ジャンプ規則を使用する。

ウィンドウ $[l, l+W_t)$ 内のカウント $c$ が $\theta$ 未満の場合、
次のイベントを含む最小の窓開始点へジャンプする：
\begin{equation}\label{eq:jump_empty}
  l \leftarrow \max(l+1,\; L_i(\bm{s}_{\mathrm{sp}})[\mathrm{next}] - W_t + 1).
\end{equation}
カウント $c > \theta$ の場合（余剰 $\mathrm{surplus} = c - \theta$ 個）、
左端から $\mathrm{surplus}$ 番目のイベントの時刻へジャンプする：
\begin{equation}\label{eq:jump_dense}
  l \leftarrow L_i(\bm{s}_{\mathrm{sp}})[l_{\mathrm{lo}} + \mathrm{surplus}].
\end{equation}
$l_{\mathrm{lo}}$ はリスト中で $l$ 以上の最初のインデックスである。

\begin{lemma}[ストライドの正確性]\label{lem:stride}
式~\eqref{eq:jump_dense} のジャンプ先 $l'$ に対し、
全ての $l'' \in (l, l')$ についても $c(l'') \geq \theta$ が成立する。
すなわち、スキップされた位置も全て密集窓の開始点である。
\end{lemma}

\begin{proof}
$l' = L_i[l_{\mathrm{lo}} + \mathrm{surplus}]$ とおく。
任意の $l'' \in (l, l')$ に対し、
ウィンドウ左端を $l$ から $l''$ に進めると
$L_i[l_{\mathrm{lo}}], \ldots, L_i[l_{\mathrm{lo}} + \mathrm{surplus} - 1]$
のうち $l''$ 未満の要素がウィンドウから脱出するが、
$L_i[l_{\mathrm{lo}} + \mathrm{surplus}] = l' > l''$ より
$L_i[l_{\mathrm{lo}} + \mathrm{surplus}]$ はまだウィンドウ内に残る。
したがって $c(l'') \geq c(l) - \mathrm{surplus} = \theta$。
\end{proof}

補題~\ref{lem:stride}は Apriori-Window の命題~3~\cite{aprioriwindow}の
$d$ 次元版であり、スキップ区間内の密集窓の開始点を
陽に評価せずとも補填により正確に列挙できることを示す。

\begin{algorithm}[t]
\caption{$t$ 軸 1 次元ストライドによる密集窓の開始点列挙}
\label{alg:t_stride}
\begin{algorithmic}[1]
\Require ソート済みイベントリスト $L$、ウィンドウ幅 $W_t$、閾値 $\theta$
\Ensure 密集窓の開始点の集合 $\mathcal{S}$
\State $\mathcal{S} \leftarrow \emptyset$,\quad
       $l \leftarrow L[0] - (W_t - 1)$,\quad
       $\mathit{prev} \leftarrow \mathrm{nil}$
\While{$l \leq L[\mathrm{last}]$}
  \State $l_{\mathrm{lo}} \leftarrow \mathrm{lower\_bound}(L, l)$,\quad
         $l_{\mathrm{hi}} \leftarrow \mathrm{upper\_bound}(L, l + W_t - 1)$
  \State $c \leftarrow l_{\mathrm{hi}} - l_{\mathrm{lo}}$
  \If{$c \geq \theta$}
    \If{$\mathit{prev} \neq \mathrm{nil}$}
      \State $\mathcal{S} \leftarrow \mathcal{S} \cup \{\mathit{prev}+1, \ldots, l-1\}$
             \Comment{ストライドで飛ばした区間を補填}
    \EndIf
    \State $\mathcal{S} \leftarrow \mathcal{S} \cup \{l\}$,\quad $\mathit{prev} \leftarrow l$
    \If{$c > \theta$}
      \State $l \leftarrow \max(l+1, L[l_{\mathrm{lo}} + (c - \theta)])$
             \Comment{式~\eqref{eq:jump_dense}}
    \Else
      \State $l \leftarrow l + 1$
    \EndIf
  \Else
    \State $\mathit{prev} \leftarrow \mathrm{nil}$
    \State $l \leftarrow \max(l+1, L[l_{\mathrm{hi}}] - W_t + 1)$
           \Comment{式~\eqref{eq:jump_empty}；$L[l_{\mathrm{hi}}]$ が次イベント}
  \EndIf
\EndWhile
\Return $\mathcal{S}$
\end{algorithmic}
\end{algorithm}

% ==============================================================
\subsection{Apriori 枝刈りによる候補セル限定（$k \geq 2$）}\label{sec:alg_k2}
% ==============================================================

$k \geq 2$ のパターン $P = \{i_1, \ldots, i_k\}$ について、
補題~\ref{lem:apriori} より
\begin{equation}\label{eq:candidate_cells}
  \mathcal{L}_\theta(S_P)
  \subseteq \bigcap_{Q \subset P,\, |Q|=k-1} \mathcal{L}_\theta(S_Q).
\end{equation}
右辺の交差集合を\emph{候補セル集合} $\mathcal{C}(P)$ と定義する。
これは Apriori-Window の「候補区間交差」~\cite{aprioriwindow}の $d$ 次元版であり、
密集セル集合の交差として定式化される。
$\mathcal{C}(P) = \emptyset$ ならば $P$ の密集セルは存在しないと確定し、
$I_P$ の構築・接頭辞和の計算を省略できる。

$\mathcal{C}(P) \neq \emptyset$ の場合は、
$\mathcal{C}(P)$ 内のトランザクションのみを走査して $I_P$ を構築し（$O(N)$）、
接頭辞和 $\mathrm{PS}_P$（$O(G)$）・サポート曲面（$O(G)$）を計算する。

% ==============================================================
\subsection{アルゴリズム全体と計算量}\label{sec:complexity}
% ==============================================================

\begin{algorithm}[t]
\caption{$d$ 次元 Apriori-Window アルゴリズム}
\label{alg:md_apriori}
\begin{algorithmic}[1]
\Require $d$ 次元トランザクション $\mathcal{D}$、パラメータ $(\bm{W}, \theta, L_{\max})$
\Ensure 各パターン $P$ の密集領域集合 $\{\mathcal{R}_P\}$
\State \textbf{離散化}: $\mathcal{D}$ を $(\bm{b}_n, d_n)$ に変換、グリッドサイズ $(N_0, \ldots, N_{d-1})$ を決定
\ForAll{$i \in \mathcal{I}$ \textbf{in parallel}}  \Comment{$k=1$; Rayon 並列}
  \State $L_i(\bm{s}_{1:}) \leftarrow $ 空間イベントリスト構築（\S\ref{sec:t_stride}）
  \ForAll{$\bm{s}_{1:} \in \prod_{k=1}^{d-1} [0, N_k)$}
    \If{$L_i(\bm{s}_{1:}) \neq \emptyset$}
      \State $\mathcal{S} \leftarrow $ \Call{TStrideStarts}{$L_i(\bm{s}_{1:}), W_0, \theta$}
             \Comment{Algorithm~\ref{alg:t_stride}}
      \State 各 $s_0 \in \mathcal{S}$ を密集セル集合 $D_i$ に追加
    \EndIf
  \EndFor
  \State $\mathcal{R}_{\{i\}} \leftarrow$ $2d$-近傍 Union-Find 連結成分 of $D_i$
\EndFor
\State $\mathcal{F}_1 \leftarrow \{\{i\} : \mathcal{R}_{\{i\}} \neq \emptyset\}$
\While{$\mathcal{F}_{k-1} \neq \emptyset$ \textbf{and} $k \leq L_{\max}$}  \Comment{$k \geq 2$}
  \State \texttt{cands} $\leftarrow$ \Call{GenerateCandidates}{$\mathcal{F}_{k-1}$}
  \State \texttt{cands} $\leftarrow$ \Call{PruneCandidates}{\texttt{cands}, $\mathcal{F}_{k-1}$}
  \ForAll{$P \in$ \texttt{cands} \textbf{in parallel}}
    \State $\mathcal{C}(P) \leftarrow \bigcap_{Q \subset P,|Q|=k-1} \mathcal{L}_\theta(S_Q)$
           \Comment{式~\eqref{eq:candidate_cells}}
    \If{$\mathcal{C}(P) = \emptyset$} \textbf{continue} \EndIf
    \State $I_P \leftarrow$ \Call{BuildItemGrid}{$\mathcal{D}, P, \mathcal{C}(P)$}
    \State $\mathrm{PS}_P \leftarrow$ \Call{BuildPrefixSumND}{$I_P$}
    \State $S_P \leftarrow$ \Call{ComputeSupportSurface}{$\mathrm{PS}_P, \bm{W}$}
    \State $\mathcal{R}_P \leftarrow$ $2d$-近傍 Union-Find 連結成分 of
           $\{\bm{s} : S_P(\bm{s}) \geq \theta\}$
  \EndFor
  \State $\mathcal{F}_k \leftarrow \{P : \mathcal{R}_P \neq \emptyset\}$,\quad
         $k \leftarrow k + 1$
\EndWhile
\Return $\{R_P : P \in \bigcup_k \mathcal{F}_k\}$
\end{algorithmic}
\end{algorithm}

\begin{table}[t]
\centering
\caption{$d$ 次元 Apriori-Window 各ステップの計算量。
$N$: トランザクション数、$G = \prod_{k=0}^{d-1} N_k$、
$W_{\mathrm{sp}} = \prod_{k=1}^{d-1} W_k$（空間ウィンドウ積）、$C_k$: $k$ 段候補数。}
\label{tab:complexity}
\small
\begin{tabular}{lll}
\toprule
ステップ & 計算量 & 備考 \\
\midrule
離散化 & $O(N)$ & \\
$k=1$: 空間イベントリスト構築 & $O(N \cdot W_{\mathrm{sp}})$ & 全アイテム合計 \\
$k=1$: $t$ 軸ストライド & $O(|\mathcal{S}_i| + \log N_i)$ per 空間キー & $|\mathcal{S}_i|$: 密集点数 \\
$k \geq 2$: 接頭辞和構築 & $O(2^d \cdot G)$ per 候補 & Apriori 枝刈り後 \\
$k \geq 2$: $I_P$ 構築 & $O(N)$ per 候補 & \\
連結成分 & $O(|\mathcal{L}| \cdot \alpha)$ per パターン & $\alpha$: 逆アッカーマン \\
\bottomrule
\end{tabular}
\end{table}

支配的コストは $k=1$ の空間イベントリスト構築 $O(N W_{\mathrm{sp}})$ と、
$k \geq 2$ の接頭辞和構築 $O(C_k \cdot 2^d \cdot G)$（$C_k$: Apriori 枝刈り後の候補数）である。
ナイーブ実装の $O\!\left(M \cdot G \cdot \prod_{k=0}^{d-1} W_k\right)$ に対して、
$t$ 軸ストライドはウィンドウ幅方向の削減、
Apriori 枝刈りはパターン次元の削減を担う。
次元数 $d$ が増えると $2^d$ と $G$ の双方が増加するが、
実用的な範囲（$d \leq 5$）では管理可能である（EX-E; \ref{sec:ex_e}節）。

% ==============================================================
\section{実験}\label{sec:experiments}
% ==============================================================

以下の研究課題を検証する。
\textbf{RQ1}: 1次元設定（空間次元なし）と比べて3次元設定は空間的偽陽性を抑制できるか？
\textbf{RQ2}: データサイズ $N$ に対して計算時間はどうスケールするか？
\textbf{RQ3}: 各高速化技術はどの程度の効果をもたらすか？（アブレーション）
\textbf{RQ4}: 既存手法（FIM、1次元 AW、Span\_3D、MZoom\_3D）と比べてどの程度優れているか？
\textbf{RQ5}: 次元数 $d$ の増加に対して計算時間はどうスケールするか？

\subsection{実験環境}

すべての実験は Rust（release build、Rayon 並列化）で実装した
\texttt{md\_pipeline} バイナリを用い、
Apple M3（8 コア、24 GB RAM）上で実行した。
デフォルトパラメータ：$\theta = 3$（サポート閾値）、
$W_t = 5$（時間ウィンドウ）、$W_x = W_y = 3$（空間ウィンドウ）、
$L_{\max} = 3$（最大パターン長）、$\Delta t = \Delta x = \Delta y = 1.0$。

\subsection{EX-A: 3次元 vs 1次元相当の比較}
\label{sec:ex_a}

\paragraph{目的.}
Apriori-Window~\cite{aprioriwindow}は時間軸のみを扱うため、
空間的に局在するパターンを識別できない（全空間を集約する）。
本実験では合成データを用いて、
3次元設定が空間的偽陽性を抑制することを実証する。

\paragraph{合成データ設計.}
$N = 5{,}000$ トランザクション、アイテム数 $|{\mathcal{I}}| = 10$、
空間範囲 $[0, 50]^2$、時刻範囲 $[0, 100]$ の合成データを生成する。
3つの密集領域を植え付ける：各領域は2アイテムパターン
$P_1 = \{1,2\}$、$P_2 = \{3,4\}$、$P_3 = \{5,6\}$ を割り当て、
それぞれランダムな時空間中心（$t_c, x_c, y_c$）の近傍に
20件の集中トランザクションを注入する。
残りの $N - 60$ 件はバックグラウンドノイズ（ランダム）。

\paragraph{比較設定.}
\begin{itemize}
\item \textbf{3次元設定}（本提案）：$W_x = W_y = 3$（植え付け領域に対応する局所ウィンドウ）
\item \textbf{1次元相当設定}：$W_x = W_y = 50$（空間範囲全体をカバーする大ウィンドウで
  空間次元を事実上無視し、1次元集約と同等の挙動を実現）
\end{itemize}
どちらも $\theta = 3$、$W_t = 5$ を使用する。

\paragraph{評価指標.}
検出された $k=2$ パターンの集合と植え付けパターン集合の一致で
Precision・Recall を計算する。5シードで実験する。

\paragraph{結果.}
表~\ref{tab:ex_a}に結果を示す。

\begin{table}[t]
\centering
\caption{EX-A: 3次元設定 vs 1次元相当設定の比較（5シード、$k=2$ パターン）。}
\label{tab:ex_a}
\small
\begin{tabular}{lrrcc}
\toprule
設定 & 検出 $k{\geq}2$ パターン数 & 植え付け回収数 & Precision & Recall \\
\midrule
3次元（$W_x{=}W_y{=}3$）     & 3   & 3 & 1.00 & 1.00 \\
1次元相当（$W_x{=}W_y{=}50$）& $\approx$93--100 & 3 & $\approx$0.03 & 1.00 \\
\bottomrule
\end{tabular}
\end{table}

\paragraph{分析.}
(RQ1)
3次元設定では全5シードで植え付け3パターン $\{1,2\}$、$\{3,4\}$、$\{5,6\}$ を
完全回収し、偽陽性パターンは0件であった（Precision $= 1.0$、Recall $= 1.0$）。
一方、1次元相当設定では $\{1,3\}$、$\{1,4\}$、$\{2,3\}$ などの
バックグラウンドノイズ由来の偽陽性パターンが多数（93〜100件）検出された
（Precision $\approx 0.03$）。

これは、空間次元を集約すると空間的に異なる位置のトランザクションが
同一ウィンドウに混入し、偶発的な共起が閾値を超えることを示す。
空間次元の明示的なモデル化が精度向上に本質的に重要である。

\subsection{EX-B: データサイズに対するスケーラビリティ}
\label{sec:ex_b}

\paragraph{実験設計.}
$N \in \{1{,}000,\; 5{,}000,\; 10{,}000,\; 50{,}000,\; 100{,}000\}$ と
ウィンドウサイズ $W_x = W_y \in \{1, 2, 3, 5, 7, 10\}$（$N = 20{,}000$ 固定）を変化させ、
実行時間を計測する。

\paragraph{結果.}
表~\ref{tab:ex_b_n}（$N$ 方向）と表~\ref{tab:ex_b_w}（ウィンドウ方向）に結果を示す。

\begin{table}[t]
\centering
\caption{EX-B: $N$ に対するスケーラビリティ（$W_x{=}W_y{=}3$、$W_t{=}5$）。}
\label{tab:ex_b_n}
\small
\begin{tabular}{rrc}
\toprule
$N$ & 検出パターン数 & 実行時間 (s) \\
\midrule
1,000   & 9  & 0.005 \\
5,000   & 13 & 0.021 \\
10,000  & 13 & 0.030 \\
50,000  & 55 & 0.455 \\
100,000 & 59 & 1.407 \\
\bottomrule
\end{tabular}
\end{table}

\begin{table}[t]
\centering
\caption{EX-B: 空間ウィンドウサイズに対するスケーラビリティ（$N{=}20{,}000$、$W_t{=}5$）。}
\label{tab:ex_b_w}
\small
\begin{tabular}{rrc}
\toprule
$W_x = W_y$ & 検出パターン数 & 実行時間 (s) \\
\midrule
1  & 13 & 0.010 \\
2  & 13 & 0.034 \\
3  & 22 & 0.086 \\
5  & 55 & 0.477 \\
7  & 55 & 1.150 \\
10 & 86 & 1.991 \\
\bottomrule
\end{tabular}
\end{table}

\paragraph{分析.}
(RQ2) 表~\ref{tab:ex_b_n}より、実行時間は $N$ に対してほぼ線形にスケールする。
$N = 100{,}000$ でも約 1.4 秒を達成しており、大規模データでの実用性を示す。

表~\ref{tab:ex_b_w}より、空間ウィンドウサイズの増大に対しては
超線形のスケーリングが観察される。
$W_x = W_y = 1$ から $W_x = W_y = 3$ まではほぼ $W^2$ に比例するが（0.010s → 0.086s、$8.6\times$）、
$W_x = W_y = 7$ 以上では空間イベントリスト構築の $O(N \prod_{k \geq 1} W_k)$ コストが支配的となる。
この結果は計算量分析（表~\ref{tab:complexity}）と整合的である。

\subsection{EX-C: 高速化技術のアブレーション}\label{sec:ex_c}

\paragraph{実験設計.}
Apriori 枝刈りの効果を測定するため、
$L_{\max}$ を 1 から 2、3、4 と増やしながら実行時間を計測する。
真の頻出パターンが $k=2$ で打ち止まる合成データ（$N = 20{,}000$）を使用する。

\paragraph{結果と分析.}
(RQ3) $L_{\max}$ を 1 から 2、3、4 と増やしても
実行時間はほぼ不変（0.062→0.083→0.082→0.087 s）であった。
これは Apriori 枝刈りが $k \geq 3$ の候補を完全に排除していることを示す。
当該データでは $k = 3$ を満たすパターンが存在せず、
$\mathcal{C}(P) = \emptyset$ の枝刈りにより $I_P$ 構築が省略されたためである。
この挙動は Apriori-Window~\cite{aprioriwindow}の1次元枝刈り効果と同様であり、
$d$ 次元への拡張でも同じ恩恵が得られることを示す。

\subsection{EX-D: 比較手法との性能比較}
\label{sec:ex_d}

\paragraph{目的.}
本実験では5種類の手法を合成データ上で比較し、
3次元 Apriori-Window の有効性を検証する（RQ4）。

\paragraph{比較手法.}
\begin{itemize}
\item \textbf{FIM\_global}: グローバル頻出アイテムセットマイニング~\cite{agrawal1994apriori}。
  空間・時間構造を無視し、全トランザクション上の共起頻度のみを使用する。
\item \textbf{1D\_AW}: Apriori-Window~\cite{aprioriwindow}（1次元、時間軸のみ）。
  空間次元を集約し、時刻スライディングウィンドウでサポートを計算する。
\item \textbf{Span\_3D}: 空間的点密度手法。
  $(1 \times 1 \times 1)$ セルごとの共起頻度が閾値を超えるパターンを検出する。
\item \textbf{MZoom\_3D}: 密部分テンソル検出~\cite{shin2018mzoom}の3次元近似。
  時空間グリッドブロック内の最大密度に基づいてパターンを検出する。
\item \textbf{3次元 AW}（本提案）: 3次元スライディングウィンドウサポートを用いた提案手法。
\end{itemize}

\paragraph{実験設計.}
$N = 5{,}000$ トランザクション（アイテム数 $|\mathcal{I}| = 20$、空間範囲 $[0, 50]^2$、
時刻範囲 $[0, 100]$）の合成データを生成し、
3つの密集領域を植え付ける（各領域 $P_1 = \{1,2\}$、$P_2 = \{3,4\}$、$P_3 = \{5,6\}$、
各20件注入、ジッタ：$t \pm 2.0$、$x \pm 1.0$、$y \pm 1.0$）。
バックグラウンドノイズは全アイテム（1--20）のランダム組み合わせ。

全手法でグリッドサーチを実施し（パラメータ範囲を表~\ref{tab:ex_d_grid}に示す）、
F1 を最大化するパラメータを選択する。
5シードで実験し平均を報告する。

\begin{table}[t]
\centering
\caption{EX-D: グリッドサーチのパラメータ範囲。}
\label{tab:ex_d_grid}
\small
\begin{tabular}{lp{6cm}}
\toprule
手法 & 探索パラメータ \\
\midrule
FIM\_global   & $\theta \in \{3,5,7,10,15,20,30\}$ \\
1D\_AW        & $W_t \in \{3,5,7,10\}$、$\theta \in \{3,5,7,10\}$ \\
Span\_3D      & $\theta \in \{3,5,7,10,15\}$ \\
MZoom\_3D     & $\theta \in \{3,5,7\}$、最小密度 $\in \{0.1,0.3,0.5\}$ \\
3次元 AW      & $W_t \in \{3,5,7,10\}$、$W_x{=}W_y \in \{2,3,5,7\}$、$\theta \in \{3,5,7,10\}$ \\
\bottomrule
\end{tabular}
\end{table}

\paragraph{結果.}
表~\ref{tab:ex_d}に各手法の最良パラメータにおける検出性能を示す。

\begin{table}[t]
\centering
\caption{EX-D: 比較実験結果（合成データ、5シード平均±標準偏差）。
植え付け $k=2$ パターン3件の回収率で評価。}
\label{tab:ex_d}
\small
\begin{tabular}{lcccc}
\toprule
手法 & Precision & Recall & F1 & 検出パターン数 \\
\midrule
FIM\_global  & $0.016 \pm 0.000$ & $\mathbf{1.000} \pm 0.000$ & $0.031 \pm 0.000$ & 190.0 \\
1D\_AW       & $0.950 \pm 0.100$ & $\mathbf{1.000} \pm 0.000$ & $0.971 \pm 0.057$ & 3.2   \\
Span\_3D     & $\mathbf{1.000} \pm 0.000$ & $0.733 \pm 0.133$ & $0.840 \pm 0.080$ & 2.2   \\
MZoom\_3D    & $\mathbf{1.000} \pm 0.000$ & $\mathbf{1.000} \pm 0.000$ & $\mathbf{1.000} \pm 0.000$ & 3.0 \\
\textbf{3次元 AW}（本提案） & $\mathbf{1.000} \pm 0.000$ & $\mathbf{1.000} \pm 0.000$ & $\mathbf{1.000} \pm 0.000$ & \textbf{3.0} \\
\bottomrule
\end{tabular}
\end{table}

\paragraph{分析.}
(RQ4)
\textbf{FIM\_global} は全3パターンを回収（Recall $= 1.0$）するが、
190件の偽陽性パターンを生成する（Precision $= 0.016$）。
これは空間・時間構造を無視することで、
バックグラウンドノイズ由来の偶発的共起が全てのパターンとして報告されるためである。

\textbf{1D\_AW}（Apriori-Window~\cite{aprioriwindow}）は時間軸のみを使用し、
大半のシードで完全回収を達成する（Recall $= 1.0$、F1平均 $= 0.971$）が、
1シード（seed=456）で偽陽性パターン1件が発生した（検出数4件）。
これは時刻が近接する異なる空間位置のトランザクションが同一時間ウィンドウに混入するためである。

\textbf{Span\_3D} は高精度（Precision $= 1.0$）だが、
平均 Recall $= 0.733$（3シードで2件回収、2シードで3件回収）。
ウィンドウ密度ではなく点密度を使用するため、
注入領域（ジッタ幅 $\pm 2$ の空間的広がり）の共起が分散した場合に閾値未達となる。

\textbf{MZoom\_3D}（近似実装）はブロック密度によるパターン検出であり、
本実験では F1 $= 1.0$ を達成した。
ただし MZoom の本来の目的はテンソル異常検出であり、
アイテムセットのスライディングウィンドウサポート意味論は持たない。

\textbf{3次元 AW}（本提案）は全シードで Precision $= 1.0$、Recall $= 1.0$ を達成し、
唯一の手法として常に正確な3パターン検出を実現した。
時空間スライディングウィンドウサポートにより、
空間的に局在した密集アイテムセット領域を偽陽性なく発見できる。

\subsection{EX-E: 次元数スケーラビリティ}\label{sec:ex_e}

\paragraph{目的.}
次元数 $d$ の増加が計算時間に与える影響を定量化する（RQ5）。

\paragraph{実験設計.}
$N = 10{,}000$ トランザクション固定、
$d \in \{2, 3, 4, 5\}$ を変化させる。
各次元のビンサイズ $\Delta_k = 2.0$、
ウィンドウサイズ $W_k = 3$（全次元統一）、
閾値 $\theta = 3$、$L_{\max} = 2$。
$t$ 軸座標範囲 $[0, 100)$、空間軸 $[0, 50)$（有効グリッドサイズ
${\approx}\,50 \times 25^{d-1}$）、
密集領域は3件植え付ける。
5シードの平均実行時間を計測する。

\paragraph{結果.}
表~\ref{tab:ex_e}に結果を示す。

\begin{table}[t]
\centering
\caption{EX-E: 次元数 $d$ に対するスケーラビリティ
（$N{=}10{,}000$、$\Delta_k{=}2$、$W_k{=}3$、$\theta{=}3$、5シード平均）。}
\label{tab:ex_e}
\small
\begin{tabular}{rrrr}
\toprule
$d$ & 有効グリッドサイズ $G$ & 実行時間 (s) & 備考 \\
\midrule
2 & ${\approx}\,50 \times 25$       & $0.07 \pm 0.00$ & 1次元相当 \\
3 & ${\approx}\,50 \times 25^2$     & $0.62 \pm 0.06$ & EX-B 設定相当 \\
4 & ${\approx}\,50 \times 25^3$     & $6.45 \pm 0.94$ & \\
5 & ${\approx}\,50 \times 25^4$     & $39.9 \pm 6.1$ & \\
\bottomrule
\end{tabular}
\end{table}

\paragraph{分析.}
(RQ5)
$d$ が1増えるごとに実行時間は約8〜10倍に増加する。
これはグリッドサイズ $G$ が空間次元 1 つあたり約 25 倍拡大し、
接頭辞和・サポートサーフェス計算が $O(G)$ であることと整合する。
$d{=}5$ では約 40 秒を要するが、
$k{\geq}2$ 候補数は Apriori 枝刈りにより実用的な範囲に抑えられている。

\subsection{考察}\label{sec:discussion}

\paragraph{密集領域の数と空間ウィンドウ.}
$W_x = W_y$ の増大に伴い検出パターン数が増加する（表~\ref{tab:ex_b_w}）。
EX-A の結果と合わせると、空間ウィンドウを大きくするほど
偽陽性パターンが増加することがわかる。
適切な空間ウィンドウサイズはデータの空間粒度に依存し、
応用ドメインの知識に基づいた設定が重要である。

% ==============================================================
\section{結論}\label{sec:conclusion}
% ==============================================================

本研究では、Apriori-Window~\cite{aprioriwindow}を
任意の $d$ 次元グリッド（時間 + $d{-}1$ 個の空間・カテゴリ次元）に拡張した
\textbf{$d$ 次元 Apriori-Window アルゴリズム}を提案した。
3つの高速化技術---$d$ 次元接頭辞和（$2^d$ コーナー包除）、$t$ 軸ストライド最適化、Apriori 枝刈り---は
いずれも Apriori-Window の対応する技術の直接拡張であり、
ナイーブ実装の $O(M \cdot G \cdot \prod W_k)$ から
$O(M \cdot G + N \cdot \prod_{k\geq1} W_k)$ へ計算量を削減した。

合成データ実験では以下を確認した：
(1) \textbf{EX-A}: 3次元設定（$W_x=W_y=3$）は植え付け3パターンを
Precision $= 1.0$、Recall $= 1.0$ で完全回収したのに対し、
1次元相当設定（$W_x=W_y=50$）は背景ノイズ由来の偽陽性パターンを
多数検出した（Precision $\approx 0.03$）。
空間次元の明示的なモデル化が精度に本質的に重要であることを示す。
(2) \textbf{EX-B}: $N = 100{,}000$ トランザクションで約 1.4 秒を達成し、
$N$ に対してほぼ線形のスケーリングを示した。
(3) \textbf{EX-C}: Apriori 枝刈りにより $k \geq 3$ の不要な計算が
完全に排除された。
(4) \textbf{EX-D}: 比較実験（5手法、5シード）では、
3次元 AW のみが Precision $= 1.0$、Recall $= 1.0$（F1 $= 1.0$）を
全シードで達成した。
FIM\_global（Precision $= 0.016$）は空間局在化不可、
1D\_AW（F1平均 $= 0.971$）は時刻近接による偽陽性が発生、
Span\_3D（F1平均 $= 0.840$）は点密度の厳格さによる見逃しが生じた。
(5) \textbf{EX-E}: 次元数スケーラビリティ実験（$d{=}2{\sim}5$）では、
$d$ が1増えるごとに実行時間は約8〜10倍増加し（$d{=}2$: 0.07秒、$d{=}5$: 40秒）、
グリッドサイズ $G \propto 25^{d-1}$ の拡大と整合した。

\paragraph{制約と限界.}
(1) 各次元はグリッド離散化を前提とし、ビン幅の選択が結果に影響する。
(2) 空間ウィンドウサイズが大きい場合（$\prod_{k\geq1} W_k \gtrsim 100$）、
空間イベントリスト構築コストが支配的となり実行時間が増大する。
(3) $d$ が大きいほど（特に $d \geq 5$）グリッドサイズ $G$ が指数的に拡大するため、
高次元設定では前処理による次元削減や粗粒度ビニングが現実的に必要となる。
(4) 本手法は密集領域の検出に特化しており、
検出された密集領域と外部イベントとの統計的関連付けは行わない。

\paragraph{今後の課題.}
(1) 実データ（小売 POS データ、疫学データ等）への応用と評価。
(2) 検出された密集領域を外部イベントに帰属させるパイプライン
（Apriori-Window~\cite{aprioriwindow}の後続研究~\cite{dsaa2025}）への接続。
(3) ストリーミング設定へのオンライン更新対応。
(4) 適応的グリッド解像度選択。


\bibliographystyle{splncs04}
\bibliography{../research/refs}

\end{document}
